\section{Experimental Evaluation}
\label{sec:experiments}

\subsection{Objective}
\label{subsec:objective}

This experimental study compares two execution styles for graph-centric tasks:
\emph{(i)} a \textbf{GGF} mode, where intermediate graphs remain inside the
SPARQL engine and are composed through graph-valued functions, and
\emph{(ii)} a \textbf{script} baseline, where a Python client orchestrates the
same logic through explicit HTTP SPARQL calls and local post-processing.

The main target of the comparison is \textbf{data shipping}. More precisely, we
measure how much intermediate information must cross the client boundary in the
GGF and script settings. Wall-clock time is also reported, but should be read
as a secondary signal because the local environment and the reranking backend
can dominate latency.

\subsection{Experimental Setup}
\label{subsec:setup}

All experiments were run on the MetaQA knowledge graph. The benchmark driver is
\texttt{xp-ggf-script/evaluate\_ggf\_vs\_script.py}. On the GGF side, the KB is
preloaded in the SPARQLLM CLI so that graph-valued functions operate directly
inside the RDFLib-backed engine. On the script side, the same KB is exposed
through a local HTTP SPARQL endpoint, and the benchmark measures the number of
HTTP requests and transferred bytes.

The run reported in this section uses the following configuration:
\begin{itemize}
\item \textbf{Dataset:} MetaQA.
\item \textbf{Benchmarked cases:} \texttt{paths},
\texttt{entity\_similarity\_compare}, and
\texttt{metaqa\_anchor\_search\_rerank\_local}.
\item \textbf{Script access mode:} HTTP SPARQL.
\item \textbf{Artificial per-call latency:} 10\,ms on the script side.
\item \textbf{Search+rereank batch size:} 5 MetaQA questions.
\item \textbf{Search index:} persistent local FAISS index over MetaQA anchor
entities.
\item \textbf{Rerank backend:} Groq, with the API key loaded from
\texttt{.env}.
\end{itemize}

\subsection{Executed Queries}
\label{subsec:executed-queries}

The three GGF benchmark cases are driven by the following SPARQL queries or
query batches.

\subsubsection*{\texttt{paths}}

\begin{verbatim}
PREFIX ggf: <http://ggf.org/>
PREFIX slm: <http://sparqllm/slm#>

SELECT ?pathRank ?stepPos ?src ?pred ?dst WHERE {
  BIND(
    ggf:PATHS(
      <http://metaqa.org/entity/avatar>,
      <http://metaqa.org/entity/titanic>,
      2,
      "both",
      10
    ) AS ?gOut
  )
  GRAPH ?gOut {
    ?root a slm:PathSet ; slm:hasPath ?path .
    ?path a slm:Path ; slm:rank ?pathRank ; slm:hasStep ?step .
    ?step a slm:PathStep ;
          slm:pos ?stepPos ;
          slm:src ?src ;
          slm:pred ?pred ;
          slm:dst ?dst .
  }
}
ORDER BY ?pathRank ?stepPos
\end{verbatim}

\subsubsection*{\texttt{entity\_similarity\_compare}}

\begin{verbatim}
PREFIX mrel: <http://metaqa.org/relation/>
PREFIX cand: <http://example.org/cand#>
PREFIX ggf:  <http://ggf.org/>

SELECT ?entity ?entityLabel ?score ?shared ?differing ?exclusive WHERE {
  VALUES ?refEntity { <http://metaqa.org/entity/avatar> }
  {
    SELECT ?entity ?entityLabel WHERE {
      ?entity mrel:release_year ?year .
      FILTER(?entity != <http://metaqa.org/entity/avatar>)
      BIND(REPLACE(STR(?entity), "^.*/", "") AS ?entityLabel)
    }
    ORDER BY ?entityLabel
    LIMIT 10
  }

  BIND(ggf:CBD(?refEntity) AS ?gRefCBD)
  BIND(ggf:SUMMARY(?gRefCBD, ?refEntity, 10, "both", "degree", 42)
       AS ?gRefSummary)

  BIND(ggf:CBD(?entity) AS ?gEntCBD)
  BIND(ggf:SUMMARY(?gEntCBD, ?entity, 10, "both", "degree", 42)
       AS ?gEntSummary)

  BIND(ggf:COMPARE-GRAPHS(?gRefSummary, ?gEntSummary, ?refEntity, ?entity)
       AS ?gCmp)

  GRAPH ?gCmp {
    ?root a cand:Comparison ;
          cand:sharedCount ?shared ;
          cand:differingCount ?differing ;
          cand:exclusiveCount ?exclusive ;
          cand:similarityScore ?score .
  }
}
ORDER BY DESC(?score) DESC(?shared) ?entityLabel
LIMIT 3
\end{verbatim}

\subsubsection*{\texttt{metaqa\_anchor\_search\_rerank\_local}}

This case evaluates a two-stage local pipeline:
\texttt{METAQA-SEARCH-FAISS(anchorLabel, 20, indexDir)} followed by
\texttt{METAQA-RERANK-LOCAL(questionText, gSearch)}. The benchmark executes one
batched query over five MetaQA questions:

\begin{verbatim}
PREFIX ggf:  <http://ggf.org/>
PREFIX cand: <http://example.org/cand#>

SELECT ?qid ?goldEntity ?entity ?label ?rank ?originalPosition
       ?searchScore ?confidence ?reason
WHERE {
  VALUES (?qid ?questionText ?anchorLabel ?goldEntity) {
    ("q1"
     "what are the languages spoken in the films directed by [Joel Zwick]"
     "Joel Zwick"
     <http://metaqa.org/entity/joel_zwick>)
    ("q2"
     "the films acted by [Sharon Tate] were released in which years"
     "Sharon Tate"
     <http://metaqa.org/entity/sharon_tate>)
    ("q3"
     "when did the films written by [Anthony Mann] release"
     "Anthony Mann"
     <http://metaqa.org/entity/anthony_mann>)
    ("q4"
     "what genres do the movies directed by [Clark Gregg] fall under"
     "Clark Gregg"
     <http://metaqa.org/entity/clark_gregg>)
    ("q5"
     "when did the movies written by [Emir Kusturica] release"
     "Emir Kusturica"
     <http://metaqa.org/entity/emir_kusturica>)
  }

  BIND(
    ggf:METAQA-SEARCH-FAISS(
      ?anchorLabel,
      20,
      "xp-ggf-script/data/metaqa_anchor_faiss"
    ) AS ?gSearch
  )
  BIND(ggf:METAQA-RERANK-LOCAL(?questionText, ?gSearch) AS ?gRank)

  GRAPH ?gRank {
    ?node a cand:RerankedEntity ;
          cand:entity ?entity ;
          cand:label ?label ;
          cand:rank ?rank ;
          cand:originalPosition ?originalPosition ;
          cand:searchScore ?searchScore ;
          cand:confidence ?confidence ;
          cand:reason ?reason .
  }
}
ORDER BY ?qid ?rank
\end{verbatim}

\subsection{Metrics}
\label{subsec:metrics}

For each case, we report:
\begin{itemize}
\item \textbf{Logical calls:} one graph-valued SPARQL evaluation on the GGF
side, versus several explicit HTTP calls on the script side.
\item \textbf{Transferred bytes:} total volume crossing the client boundary.
\item \textbf{Wall time:} end-to-end execution time.
\item \textbf{Functional equivalence:} whether the normalized final outputs are
identical.
\end{itemize}

\subsection{Results for the 10\,ms HTTP-Latency Run}
\label{subsec:results-10ms}

Table~\ref{tab:three-case-benchmark-results} summarizes the benchmark run with
10\,ms artificial HTTP latency. The results confirm the expected pattern:
GGF always uses a single logical call, while the script baseline issues
multiple HTTP requests and transfers substantially more data.

The strongest communication gap is observed on \texttt{paths}: the script
baseline transfers about \(70.30\times\) more bytes than GGF. For
\texttt{entity\_similarity\_compare}, the gap is still \(20.74\times\), and
for the local search+rereank pipeline it remains \(13.37\times\). The latter
case is not functionally equivalent because the reranking backend is
nondeterministic and is invoked independently in the GGF and script runs.
However, this does not invalidate the communication measurement itself, since
the benchmark still captures the fact that the script baseline must externalize
the candidate list and reranking exchanges, whereas GGF keeps the intermediate
candidate graph inside the engine.

\begin{table}[t]
\centering
\small
\begin{tabular}{l|r|r|r|r|r|c}
\hline
Task & GGF bytes & Script bytes & Ratio & GGF calls & Script calls & Eq. \\
\hline
\texttt{paths} & 3,204 & 225,241 & 70.30$\times$ & 1 & 4 & yes \\
\texttt{entity\_similarity\_compare} & 1,466 & 30,410 & 20.74$\times$ & 1 & 12 & yes \\
\texttt{metaqa\_anchor\_search\_rerank\_local} & 7,746 & 103,599 & 13.37$\times$ & 1 & 10 & no \\
\hline
\end{tabular}
\caption{Three-case benchmark on MetaQA with HTTP-based script execution and
10\,ms artificial latency per script-side SPARQL request.}
\label{tab:three-case-benchmark-results}
\end{table}

Figure~\ref{fig:three-case-benchmark-plot} gives the corresponding visual
summary. The first three panels use logarithmic scales for transferred bytes,
logical calls, and transfer ratios. The wall-time panel is kept on a linear
scale so that close values such as \(12.70\)s versus \(12.21\)s remain visually
interpretable.

\begin{figure}[t]
\centering
\includegraphics[width=\linewidth]{xp-ggf-script/out/bench_paths_similarity_search_10ms.png}
\caption{Benchmark plot for the three-case MetaQA HTTP run with 10\,ms
artificial latency. The panels show transferred bytes, logical calls, transfer
ratio, and wall-clock time.}
\label{fig:three-case-benchmark-plot}
\end{figure}

\subsection{Takeaway}
\label{subsec:takeaway}

This run reinforces the main claim of the GGF approach: \emph{intermediate
graphs do not need to be externalized}. Even on a local HTTP setup with only
10\,ms of artificial latency, the script baseline already pays a substantial
communication penalty, because it must materialize, transfer, and post-process
intermediate data explicitly. GGF collapses these steps into one server-side
graph computation and therefore reduces the amount of client-visible data
shipping by one to two orders of magnitude on the three tasks reported here.
